\section{Choose a lane width for the arithmetic}\label{rgb-lanes}
An MSA register holds 128 bits. The suffix on an instruction selects how those bits
are divided into elements: sixteen bytes, eight halfwords, four words or two
doublewords. The register does not acquire a permanent type when loaded. A byte
shuffle can arrange bytes so that a later halfword shift sees the intended values.

For packing, construct eight halfwords containing one channel each. Every halfword
must have a component byte in its low position and zero in its high position. A
halfword right shift can then discard the component's low bits without borrowing
bits from a neighboring pixel. Red and green left shifts place the fields within
each halfword; OR combines the three vectors lane by lane.

\texttt{ld.b} and \texttt{st.b} transfer sixteen bytes each. The suffix describes
byte elements, not a transfer of one byte. In the pinned little-endian environment,
loading byte elements places the first memory byte in byte lane zero. Two adjacent
byte lanes then form a little-endian halfword when the same bits are used by a
halfword operation. Keep that relationship separate from the host machine's own
byte order.

The scalar tail handles fewer than eight remaining pixels. For counts eight or more,
the vector loop handles one complete eight-pixel block at a time and updates count
by eight. Zero and short counts reach the scalar bound check before any source load.
The vector and scalar paths share the same original-count return convention.
